\section{Related Work}
%\textbf{Activity Recognition with RGB Data:} \textcite{tran2015learning} propose the ”Two-Stream Convolutional Networks” that integrated spatial and temporal streams to capture static appearance and motion information, respectively. This work demonstrated the efficacy of fusing RGB data with optical flow information for activity recognition, achieving state-of-the-art performance on benchmark datasets. The introduction of C3D by \textcite{ji20123d} leverages 3D convolutions to capture spatio-temporal patterns directly from RGB videos, demonstrating superior performance in activity recognition compared to traditional 2D CNNs. Since the introduction of I3D \cite{carreira2017quo}, 3D-CNN has emerged as the standard method for activity recognition. Since then, the activity recognition community has put out several cutting-edge 3DCNN designs \cite{feichtenhofer2020x3d, feichtenhofer2019slowfast} that exceed I3D in accuracy and efficiency.

%Moreover, \textcite{franco2020multimodal} and \textcite{tran2015learning} propose a novel activity encoding method using temporal images from RGB video sequences and integrating complementary information from a skeleton and temporal data. The proposed method outperforms state-of-the-art methods on public datasets. The challenges addressed in this paper include the difficulty of recognizing complex activities with high intra-class variability and the need for robustness to occlusions and noise. These challenges were addressed by using a multimodal approach that combines RGB and skeleton data, as well as a novel activity encoding method that uses temporal images.

\textbf{Skeleton-based activity recognition:} Skeleton data offers a compact and informative representation of human actions, focusing on the spatial arrangement and temporal dynamics of joints.
%Early approaches involved the use of histograms of joint angles or velocities as feature representations~\cite{xia2012view}. By quantizing joint angles or velocities into bins and constructing histograms, these methods attempted to capture the spatial and temporal dynamics of actions. 
With the arrival of deep learning, there was a paradigm shift in skeleton-based activity recognition. Recurrent neural networks (RNN) were amongst the first deep learning methods applied to skeleton based activity recognition \cite{baccouche2011sequential}. RNNs enabled the modeling of temporal dependencies in activity sequences, capturing the dynamics of actions over time \cite{song2017endtoend, wang2017modeling}. Convolutional neural networks (CNN) were adapted to handle skeleton data by treating joint coordinates as image-like representations \cite{du2015skeleton, zhang2019view}. \textcite{zhang2018adding} propose a recursive connection to create a EleAtt-RNN block. They address the challenges of modeling complex sequential information with fixed-weight RNNs and the limitation of existing attention mechanisms that require additional parameters and computational overhead. However, various related methods \cite{wu2014leveraging, zhao2017two} often are unable to achieve a competitive outcome due to the RNN-based architecture’s poor spatial modeling capability.

CNN architectures can effectively handle spatial data, leading to an increased interest in CNNs for skeleton based activity recognition tasks. \textcite{ding2017investigation} investigate encoding richer spatial features into texture color images or 2D pseudo-images for 3D human activity recognition using a CNN-based approach. Each type of feature is encoded into images in two or more ways to further explore the spatio-temporal information. \textcite{liu2017enhanced} investigates methods to enhance skeleton visualization for improved activity recognition under varying viewpoints. On the other hand, GCNs have emerged as a popular choice for effectively representing and processing graph-structured data, including skeleton graphs \cite{shi2019two, cheng2020skeleton, martin_dynamic_2020}. The ST-GCN approach constructs a spatial-temporal graph using joints as graph vertices and incorporating natural connections in both human body structures and time as graph edges.
Since the introduction of ST-GCN, the use of GCNs for skeleton-based activity recognition has garnered significant attention, leading to various related works. \textcite{li2019actional} propose AS-GCN, which does not only recognize a person’s activity but also employing a multi-task learning strategy to predict the subject’s next possible pose.

In skeleton-based driver activity recognition earlier work focused on RNN-based~\cite{drive_and_act_2019_iccv} and GCN-based approaches~\cite{martin_dynamic_2020}. More recent work incorporates attention mechanism with GCN for efficient driver activity recognition~\cite{li2024incomplete}.

The vanilla transformer model~\cite{vaswani2017attention} makes significant strides in natural language processing (NLP) and has since been adopted in various domains. Inspired by the self-attention in computer vision multiple approaches adopted the transformer principle for skeleton-based activity recognition~\cite{qiu2022spatio, wu2022skeletonmae, cheng_hierarchical_2021, kim_global-local_2022}. \textcite{plizzari2021skeleton} introduce a self-attention based approach instead of regular graph convolutions in both spatial and temporal dimensions. 
This replacement of the graph structure of skeleton data by self-attention allows the model to effectively capture dependencies between joints in space and time, leading to improved activity recognition performance.

\textcite{qiu2022spatio} introduce a spatio-temporal tuples transformer (STTFormer). The skeleton sequence is split into multiple patches, and each patch has numerous consecutive frames that are encoded. The link between various joints’ non-consecutive frames is then captured by a self attention module based on spatio-temporal tuples. 
%This model has shown better performance on benchmark datasets when compared to state-of-the-art methods. 

\textbf{Self-supervised skeleton-based activity recognition:} To our knowledge there is no related 3D skeleton-based driver activity recognition method using self-supervised learning. Therefore we focus this related work on self-supervised skeleton-based human activity recognition. Various self-supervised techniques have been developed to leverage the temporal dynamics and spatial information present in skeleton data. \textcite{zheng2018unsupervised} propose a method based on a recurrent encoder-decoder generative adversarial network (GAN) to reconstruct the input skeleton sequence, capturing long-term temporal dynamics and subtle motion patterns. \textcite{su2020predict} introduce an autoencoder for encoding and reconstructing the skeleton sequence, leveraging clustering and nearest neighbor classification for downstream tasks. Existing 3D skeleton-based autoencoders utilize RNN~\cite{su2020predict, wu2022bi, rao_augmented_2021} or transformer backbones~\cite{Lerch_2024_WACV, wu2022skeletonmae}.

\textcite{lin2020ms2l} proposes a multi-task framework (MS2L) for enhancing the understanding of motion and spatial relationships among joints. AS-CAL~\cite{rao_augmented_2021} introduces a momentum long short-term memory (LSTM) based approach to regularize the feature space, while SkeletonCLR~\cite{li20213d} applies a memory bank and cross-view knowledge mining for more comprehensive representations. \textcite{guo2022contrastive} propose AimCLR, relying on powerful data augmentation, and \textcite{zhou2023self} focus on leveraging local relationships between skeleton joints and video frames in their partial spatio-temporal learning (PSTL) approach. This efficient masking strategy is crucial for real-world applications. To the best of our knowledge we are the first to explore self-supervised learning for skeleton-based driver activity recognition. Inspired by the DINO framework~\cite{caron2021emerging}, we propose an adaptation explicitly designed for 3D skeleton data. Our baseline to our proposed method is the PSTL approach \cite{zhou2023self}. PSTL is tested in skeleton-based general human activity recognition.

%\textbf{Self-supervised skeleton-based activity recognition:} Besides the supervised approaches various self-supervised techniques are developed to leverage the temporal dynamics and spatial information present in skeleton data. \textcite{zheng2018unsupervised} propose a method based on a recurrent encoder-decoder GAN to reconstruct the input skeleton sequence. With learning to reconstruct the input sequence, the model can capture long-term temporal dynamics and subtle motion patterns essential for activity recognition. \textcite{su2020predict} introduce a autoencoder to encode and reconstruct the skeleton sequence. For the classification downstream task they use clustering and nearest neighbor classifier. The method leverages both prediction and clustering to encourage the model to learn more discriminative and compact representations. 
%Existing 3D skeleton-based autoencoders are based on RNN \cite{su2020predict, wu2022bi, rao_augmented_2021} or transformer \textcite{Lerch_2024_WACV, wu2022skeletonmae} backbones.
%\textcite{lin2020ms2l} propose MS2L a multi-task framework, including motion prediction and jigsaw puzzle tasks, to enhance the model’s understanding of motion and spatial relationships among joints. \textcite{rao_augmented_2021} introduce AS-CAL, an momentum LSTM based approach to regularize the feature space. \textcite{li20213d} propose SkeletonCLR, which applies a memory bank to store negative samples and employs a cross-view knowledge mining strategy. By leveraging the memory bank and cross-view consistency, the model can capture more comprehensive and discriminative representations. In order to add movement patterns and compel the encoder to acquire broader representations, \textcite{guo2022contrastive} propose AimCLR to make extensive use of augmentation. These contrastive learning approaches significantly rely on powerful data augmentation procedures. 
%While current approaches focus on strong augmentation techniques to discriminate different skeletons, \textcite{zhou2023self} aim to leverage the local relationship between various skeleton joints and video frames in their Partial Spatio-Temporal Learning (PSTL) approach. In the PSTL approach they adopt a unique spatio-temporal masking strategy to construct partial skeleton sequences, allowing the model to focus on the specific regions of interest. The framework utilizes a negative-sample-free triplet stream structure adopted from Barlow Twins (BT) \cite{zbontar2021barlow} for a so-called SkeletonBT, comprising an anchor stream without any masking, a spatial masking stream with Central Spatial Masking (CSM), and a temporal masking stream with Motion Attention Temporal Masking (MATM). This efficient masking strategy is crucial for real-world applications.

%Although this approach works well for general 3D skeleton-based human activity recognition datasets, the performance drops when applied on the car interior domain. Therefore, we propose an approach inspired by the DINO framework \cite{caron2021emerging}. Unlike the original DINO, which focuses on image data, our adaptation is designed explicitly for 3D skeleton data.

%\pagebreak